\documentclass[11pt]{beamer} %Add 'handout' in the options for printing

\usepackage{beamercolorthemewolverine}
\usetheme{AnnArbor}

\usepackage[english]{babel}
\usepackage{amsmath,amssymb}
\usepackage{mathtools}
\usepackage{commath}
\usepackage{amssymb,amscd,amsfonts,amsbsy}
\usepackage{enumerate}
\usepackage{epsfig}
\usepackage{pgfpages}
\usepackage{graphics}
\usepackage{graphicx}
\usepackage{multicol}
\usepackage{layouts}
\usepackage[center]{caption}
\usepackage{natbib}
%\usepackage{movie15}
\usepackage{etex}
\usepackage[export]{adjustbox}
\usepackage{embedfile}
\usepackage{caption}
\usepackage{soul,xcolor}
\usepackage{import}
\usepackage{relsize}
\usepackage{tikz}
\usepackage[makeroom]{cancel}
\usepackage{multimedia}
\usepackage{subcaption}
\usepackage{comment}
%\usepackage{media9}



%Setup note page
\setbeamertemplate{note page}[plain]
%\setbeameroption{show notes on second screen=right}
%\setbeameroption{show notes on second screen=left}

\setbeamertemplate{navigation symbols}{}%remove navigation symbols

\setbeamertemplate{footline}[frame number]{}%gets rid of bottom navigation bars


\graphicspath{ {../figs/lung_figs/} }


% Figures within a column...
\makeatletter
\newenvironment{tablehere}
{\def\@captype{table}}
{}
\newenvironment{figurehere}
{\def\@captype{figure}}
{}
\makeatother

\newcommand{\beginbackup}{
   \newcounter{framenumbervorappendix}
   \setcounter{framenumbervorappendix}{\value{framenumber}}
}
\newcommand{\backupend}{
   \addtocounter{framenumbervorappendix}{-\value{framenumber}}
   \addtocounter{framenumber}{\value{framenumbervorappendix}} 
}

%---------------------------------------------------------------------------
%\setbeamercolor{mycolor}{fg=White,bg=White}
\setbeamercolor{alerted text}{fg=red} 
\setbeamertemplate{frametitle}
{
  \vskip-5pt 
  \leavevmode
  \begin{center}
    \vspace*{-0.5cm}
  \hbox{%
  \begin{beamercolorbox}[wd=0.85\paperwidth,ht=0.35cm,dp=1.35ex]{White}%
    \raggedright
    \vspace*{-.25cm}
    {\normalsize \insertframetitle} %\small{\insertframetitle}
  \end{beamercolorbox}
  % \hspace*{1em}{\normalsize \insertframetitle} %\small{\insertframetitle}
  }%
  \end{center}
  \vspace*{-0.75cm}
}
\def\newblock{\hskip .11em plus .33em minus .07em}

\setbeamertemplate{headline}{}

\setbeamercovered{dynamic}


\newcommand{\RR}{{\mathbb{R}}}
\newcommand{\NN}{{\mathbb{N}}}
\newcommand{\ZZ}{{\mathbb{Z}}}
\newcommand{\CC}{{\mathbb{C}}}
\newcommand{\eps}{\varepsilon}
\newcommand{\bp}{\noindent {\it Proof}.\,\,\,}
\newcommand{\ep}{\hfill$\Box$ \vskip 0.08in}
\newcommand{\dint}{\int\!\!\!\int}
\newcommand{\vs}{\vskip 0.5cm}
\newcommand{\po}{{\partial\Omega}}
\newcommand{\meanint}{{\int{\mkern-19mu}-}}
\def\ring{\mathaccent"0017 }

\definecolor{AliceBlue}{RGB}{240,248,255}


%\definecolor{lightblue}{rgb}{1,0,0}
%\sethlcolor{lightblue}
\renewcommand<>{\hl}[1]{\only#2{\beameroriginal{\hl}}{#1}}

% http://tex.stackexchange.com/questions/41683/why-is-it-that-coloring-in-soul-in-beamer-is-not-visible
\makeatletter
\newcommand\SoulColor{%
  \let\set@color\beamerorig@set@color
  \let\reset@color\beamerorig@reset@color}
\makeatother
\SoulColor



%\usepackage{media9}%
%\newcommand{\includemovies}[3]{%
%\includemedia[%
%width=#1,height=#2,%
%activate=pagevisible,%
%deactivate=pageclose,%
%addresource=#3,%
%flashvars={%
%src=#3 % same path as in addresource!
%&autoPlay=true % default: false; if =true, automatically starts playback after activation (see option ‘activation)’
%&loop=true % if loop=true, media is played in a loop
%&controlBarAutoHideTimeout=0 %  time span before auto-hide
%}%
%]{}{StrobeMediaPlayback.swf}%
%}



%%% Local Variables:
%%% mode: latex
%%% TeX-master: "../main"
%%% End:

\embedfile{\jobname}
\embedfile{include/preamble.tex}
%

% My commands
\newcommand{\orderof}[1]{\ensuremath{\mathcal{O}\left(#1\right)}}

\begin{document}
%% Some customizations, based on https://tex.stackexchange.com/questions/22346/how-to-customize-titlepage-in-beamer

% \defbeamertemplate*{title page}{customized}[1][]
% {
% %  \usebeamerfont{title}\inserttitle\par
% %  \usebeamerfont{subtitle}\usebeamercolor[fg]{subtitle}\insertsubtitle\par
% %  \bigskip
%  % \usebeamerfont{author}\insertauthor\par
% %  \usebeamerfont{institute}\insertinstitute\par
% %  \usebeamerfont{date}\insertdate\par
%   \usebeamercolor[fg]{titlegraphic}\inserttitlegraphic
% }


%%
\title[]{Applications of Computation in Acoustics:\\Ultrasound Bioeffects \&\\ Underwater Transmission Loss Uncertainty}
\author[] {}

\institute[]{}
\date[date]{}


\begin{frame} %\vspace{1cm}
  \vspace*{1cm}
  %\titlepage \vspace{-2.50cm}

  {
    %\definecolor{darkpowderblue}{rgb}{0.0, 0.2, 0.6} %% UNCOMMENT IF NOT IN PREAMPLE
    \setbeamercolor{title}{fg=white, bg=darkpowderblue}
  \begin{beamercolorbox}[sep=8pt,center,wd=\paperwidth,ht=0.3\textheight,dp=0.045\textheight]{title}
    \usebeamerfont{title}\inserttitle\par%
  \end{beamercolorbox}
  }
  \centering
  \begin{columns}
    \column{.25\textwidth}
    
    % \begin{figure}\vspace*{1.3cm} \centering \includegraphics[width=.70\columnwidth]{include/figs/MichiganSeal} \end{figure}
    \column{.5\textwidth}
    \begin{center}
      % 
      \normalsize{A \emph{dissertation defense} by:\\Brandon Patterson} \vspace{16pt} \\ 
      % 
      \footnotesize{Scientific Computing and Flow Physics Lab} \vspace{4pt}\\
      \footnotesize{University of Michigan, Ann Arbor} \vspace{4pt}\\
      \scriptsize{Department of Mechanical Engineering} \vspace{4pt}\\
      % \scriptsize{$^2$Department of Radiology} \vspace{4pt}
    \end{center}
    \column{.25\textwidth}
    \vspace*{.15cm}
    % \begin{figure} \vspace*{1.3cm} \centering \includegraphics[width=.70\columnwidth]{include/figs/MElogo} \end{figure}
  \end{columns}
  \vspace{.25cm}
  \begin{center}
    % 
    % 
    %%%%%%%%%%%%%%%%%%% Put Date Here   %%%%%%%%%%%%%%%%%%%%%%%%%%%%
    % \footnotesize {\today} \vspace{-.1cm}
    \footnotesize {October 31, 2017} \vspace{-.1cm}
    % 
    % 
  \end{center}



  % 
  % 
  %%%%%%%%%%%%%%%%%%% FUNDING SOURCE OR CONFERENCE IMAGES BELOW ( L | C | R )   %%%%%%%%%%%%%%%%%%%%
  \begin{minipage}{\textwidth}
    \hspace*{.5cm}
    \begin{minipage}{.3\textwidth}
      \begin{center} 
        \begin{figure}
          \centering
          % \includegraphics[width=.4\textwidth]{include/figs/NIH_logo}
          %	\label{fig:NIH_logo_blue}
        \end{figure}
      \end{center}
    \end{minipage}
    \begin{minipage}{.3\textwidth}
      \begin{center} 
        \begin{figure}
          \centering
          % \includegraphics[width=.8\textwidth]{include/figs/nsfgrfp}
          \label{fig:NSF_logo_blue}
        \end{figure}
      \end{center}
    \end{minipage}
    \begin{minipage}{.3\textwidth}
      \begin{center} 
        
      \end{center}
    \end{minipage}
  \end{minipage}
\end{frame}
%%% Local Variables:
%%% mode: latex
%%% TeX-master: "../main"
%%% End:




%%%%%%%%%%%%%%%%%%%%%%%%%%%%%%%% Slides %%%%%%%%%%%%%%%%%%%%%%%%%%%%%%
\begin{frame}
  \centering
  \begin{center}
    \LARGE Part I: Introduction
  \end{center}
\end{frame}
%%
%%
\begin{frame} \frametitle{Introduction}

\end{frame}
%%
%%
\begin{frame} \frametitle{Past Project: Efficient estimation of the Probability Density Function of transmission loss in uncertain ocean environments}

\end{frame}
%%
%%
\begin{frame} \frametitle{Past Project: Theoretical microbubble dynamics in a viscoelastic medium at capillary breaching thresholds}

\end{frame}
%%
%%
\begin{frame} \frametitle{Background on diagnostic ultrasound}
  \begin{itemize}
    \item Diagnostic Ultrasound (DUS) is ubiquitous and generally very safe
    \item Lung Hemorrhage (LH) is the only known bioeffect of non-contrast diagnostic ultrasound
  \end{itemize}
\end{frame}
%%
%%
\begin{frame} \frametitle{Background on wave-interface interaction}

\end{frame}
%%
%%
\begin{frame}
  \centering
  \begin{center}
    {\LARGE Part II: Current Project}\\
    Diagnostic ultrasound-induced lung hemorrhage and acoustic wave
    interactions with liquid-gas interfaces
  \end{center}
\end{frame}
%%
%%
%%
%%
\begin{frame} \frametitle{Diagnostic ultrasound-induced lung hemorrhage is not a new problem}
  \begin{itemize}
  \item Discovered in mice during lithotripsy experiment by \cite{Child1990}.
  \item Since has been shown to occur in a variety of mammals \SoulColor{\hl{ADD CITATION}}
  \end{itemize}
\end{frame}
%%
%%
\begin{frame} \frametitle{Mechanical wave fluid-fluid interface interactions}
  \begin{itemize}
  \item 
  \end{itemize}
\end{frame}
%%
%%
\begin{frame}
We aim to study acoustic wave interactions with perturbed liquid-gas interfaces
\end{frame}
%%
%%
\begin{frame} \frametitle{Problem setup}
\end{frame}
%%
%%
\begin{frame} \frametitle{Methods}
\end{frame}
%%
%%
\begin{frame} \frametitle{Results: Behavior of the interface}
\end{frame}
%%
%%
\begin{frame} \frametitle{Results: Evolution of the interface}
  \begin{itemize}
  \item Plot of interface amplitude vs time
  \item Scaling law
  \end{itemize}
\end{frame}
%%
%%
\begin{frame} \frametitle{Results: Vorticity dynamics}
  \begin{itemize}
  \item Surface plots of vorticity throughout the interaction
  \item Two questions: What kinds of vorticity, can we predict where it is generated
  \end{itemize}
\end{frame}
%%
%%
\begin{frame} \frametitle{Analysis of vorticity dynamics}%
  \begin{itemize}%
  \item Analysis of vorticity dynamics
  \end{itemize}
\end{frame}
%%
%%
\begin{frame} \frametitle{Results: Vorticity dynamics}
  \begin{itemize}
  \item Surface plot of vorticity during the compression wave interaction
  \item Plot of vorticity vs time (during wave interaction)
  \item Vorticity distribution plots
  \end{itemize}
\end{frame}
%%
%%
\begin{frame} \frametitle{Implications for DUS}
  \begin{itemize}
  \item 
  \end{itemize}
\end{frame}
%%
%%
\begin{frame} \frametitle{Conclusions thus far}

\end{frame}
%%
%%

\begin{frame} \frametitle{Vorticity generation vs $y_0$ across interface}
%See circ_calcs.m (attached)
{\scriptsize
  \begin{align}%
    \label{eq:baroclinic_air_water}%
    \frac{\norm{\frac{\nabla\rho\times\nabla p}{\rho^2}}_{air\quad}}{\norm{\frac{\nabla\rho\times\nabla p}{\rho^2}}_{water}}%
    =&\orderof{\abs{\boldsymbol{T}}\left(\frac{\abs{\rho^-}}{\abs{\rho^+}}\right)^2}%
    \approx 357
  \end{align}
  \vspace*{-.5cm}
  %
  \begin{figure}
    \hfill
    \includegraphics[width=.45\textwidth]{./figs/vorticity_vs_y0} \hfill
    \includegraphics[width=.45\textwidth]{./figs/circ_y0_dist}
    \hfill
  \end{figure}
  %
  \vspace*{-.5cm}
  \begin{itemize}
    \item 97\% of circulation appears in fluid with $y_0<0.5$
    \item Computed circulation ratio from air($y_0<0.1$) to water($y_0>0.9$) is $\approx 27$
  \end{itemize}
}
\end{frame}
%%
%%
\begin{frame}
  \centering
  \begin{center}
    \LARGE Part III: Future work
  \end{center}
\end{frame}
%%%%%%%%%%%%%%%%%%%%%%%%%%%%%%%%%%%%%%%%%%%%%%%%%%%%%%%%%%%%%%%%%%%%%% 


\section{Citations}
\bibliographystyle{../tex/myauthordate2}
% Give this command the relative path to the .bib file.
\bibliography{../../../literature/library}


%%%%%%%%%%%%%%%%%%%%%%%%%%%%% 

% \begin{frame} \vspace{\fill} \begin{center} \Huge 
%     Thanks a lot!\\ \vspace{\fill} Questions? 
%   \end{center} \vspace{\fill} \end{frame} 


\end{document}






% Template Created by Brandon Patterson - awesome@umich.edu
%%% Local Variables:
%%% mode: latex
%%% TeX-master: t
%%% End:
